\chapter{Fibrinogen}

\section*{What Is This Test?}
Fibrinogen (factor I) is a 340-kDa soluble glycoprotein synthesized by the liver that circulates in plasma at concentrations of 200--400 mg/dL. It is the central protein of the coagulation cascade: thrombin cleaves fibrinogen to form fibrin monomers, which polymerize and are cross-linked by factor XIIIa to form the insoluble fibrin clot that is the structural framework of hemostasis. Fibrinogen is also an acute-phase reactant whose concentration rises 2--10 fold in response to inflammation, infection, tissue injury, and malignancy through interleukin-6 (IL-6)-stimulated hepatic synthesis. The fibrinogen assay quantifies the concentration of functional (clottable) fibrinogen in plasma and is used to diagnose congenital fibrinogen disorders (afibrinogenemia, hypofibrinogenemia, dysfibrinogenemia), to evaluate and monitor disseminated intravascular coagulation (DIC), to assess liver synthetic function, to investigate unexplained bleeding or thrombosis, and to monitor fibrinogen replacement therapy in massive hemorrhage. A low fibrinogen is a key component of the diagnostic criteria for DIC. A high fibrinogen is an independent risk factor for cardiovascular disease (atherosclerosis, myocardial infarction, stroke) and venous thromboembolism.

\section*{Alternative Names}
Factor I; Fibrinogen Activity; Functional Fibrinogen; Clauss Fibrinogen; Fibrinogen Antigen; Plasma Fibrinogen

\section*{Principle}
Fibrinogen is most commonly measured by the Clauss method, a modified thrombin time. Citrated platelet-poor plasma is diluted 1:10 with buffer, and a high concentration of bovine or human thrombin (approximately 100 NIH units/mL) is added. The thrombin converts fibrinogen to fibrin, and the clotting time is measured. Because the plasma is diluted and excess thrombin is used, the clotting time is inversely proportional to the fibrinogen concentration (inhibitors such as heparin and fibrin degradation products are diluted out, minimizing interference). The fibrinogen concentration is determined by comparing the clotting time to a standard curve prepared from a reference plasma with known fibrinogen concentration. Results are reported in mg/dL or g/L. Alternative methods include: the PT-derived fibrinogen method (estimated from the change in optical density during PT clot formation --- less accurate than Clauss, especially at low fibrinogen levels); immunologic methods (ELISA, immunoturbidimetric) that measure fibrinogen antigen regardless of function (used to detect dysfibrinogenemia where functional fibrinogen is low but antigen is normal); and the heat precipitation method of Schulz (rarely used).

\section*{History}
Fibrinogen was first identified as the precursor of fibrin in the early 19th century. The term ``fibrinogen'' was coined by Rudolf Virchow in 1847. The first quantitative method for fibrinogen was developed by Olof Hammarsten in 1879. The Clauss method was developed by Alfons Clauss in 1957 and remains the gold standard for functional fibrinogen measurement. The recognition of fibrinogen as an acute-phase reactant and an independent cardiovascular risk factor emerged in epidemiological studies in the 1980s and 1990s, with the Northwick Park Heart Study (1986) demonstrating that elevated fibrinogen was as strong a predictor of coronary heart disease as elevated cholesterol. The identification of congenital fibrinogen disorders (afibrinogenemia, hypofibrinogenemia, dysfibrinogenemia) as distinct clinical entities progressed through the 20th century.

\section*{Why This Test Is Ordered}
\begin{itemize}
  \item Diagnosis and monitoring of disseminated intravascular coagulation (DIC) --- low or decreasing fibrinogen is a key criterion in the ISTH DIC scoring system; marked consumption leads to hypofibrinogenemia
  \item Evaluation of congenital fibrinogen disorders --- afibrinogenemia (absent, severe bleeding in infancy), hypofibrinogenemia (low, variable bleeding), dysfibrinogenemia (normal antigen but abnormal function; variable: bleeding, thrombosis, or asymptomatic)
  \item Identification of massive hemorrhage with dilutional and consumptive coagulopathy --- fibrinogen is the first coagulation factor to reach critically low levels during massive transfusion
  \item Assessment of liver synthetic function --- fibrinogen is synthesized exclusively in the liver; low levels in advanced cirrhosis and acute liver failure
  \item Investigation of unexplained bleeding or prolonged PT/aPTT/thrombin time
  \item Preoperative screening in patients with bleeding history
  \item Monitoring of fibrinogen replacement therapy (cryoprecipitate, fibrinogen concentrate) in massive hemorrhage
  \item Cardiovascular risk assessment --- elevated fibrinogen is an independent predictor of coronary artery disease, myocardial infarction, and stroke
  \item Monitoring of acute-phase response in inflammatory conditions
\end{itemize}

\section*{Is This a Primary or Secondary Test}
Fibrinogen measurement is a secondary coagulation test ordered when PT and aPTT are prolonged, when DIC is suspected, when congenital fibrinogen disorders are considered, or when cardiovascular risk stratification is desired.

\section*{How This Test Is Performed}
A venous blood sample is collected in a light blue-top tube containing 3.2\% sodium citrate (9:1 ratio), filled exactly to the line. The sample is centrifuged at 1,500--2,500 $\times$ g for 15 minutes to obtain platelet-poor plasma. For the Clauss method, plasma is diluted 1:10 with buffer, and a high concentration of thrombin is added. The clotting time is measured by optical or mechanical detection on an automated coagulation analyzer. The fibrinogen concentration is calculated from a standard curve. The venipuncture takes less than 5 minutes. Analysis is completed in 2--5 minutes.

\section*{What Preparation Is Needed}
No special preparation is required. Fasting is not necessary. The patient should disclose all medications, especially heparin, warfarin, direct oral anticoagulants, thrombolytic agents (streptokinase, tPA), fibrinogen concentrates, cryoprecipitate, and recent blood transfusion. Acute infection, inflammation, pregnancy, and recent surgery/trauma should be disclosed (all elevate fibrinogen as an acute-phase reactant).

\section*{Contraindications and Precautions}
No absolute contraindications. Pre-analytical factors: the light blue-top tube must be filled exactly to the line (9:1 ratio); hemolyzed, clotted, or lipemic specimens are unacceptable; traumatic venipuncture activates coagulation and consumes fibrinogen; heparin contamination from lines falsely prolongs clotting time and underestimates fibrinogen; heparin effect is minimized in the Clauss method by the 1:10 dilution but very high heparin levels may still interfere; elevated levels of fibrin degradation products (FDPs) in DIC inhibit fibrin polymerization and cause falsely low Clauss fibrinogen results; fibrinogen is an acute-phase reactant and interpretation must account for clinical context (a ``normal'' fibrinogen in the setting of severe sepsis/DIC may actually represent significant consumption from an initial high baseline); fibrinogen levels decline with prolonged storage at room temperature; testing should be performed within 4 hours of collection at room temperature, or plasma frozen at -20\textdegree{}C.

\section*{Risks}
Minimal risks of venipuncture: mild pain, bruising, hematoma, lightheadedness, and rarely infection or nerve injury.

\section*{What Happens After the Test}
Results are available within 30--60 minutes. In DIC evaluation, fibrinogen is interpreted with platelet count, PT, aPTT, and D-dimer. Fibrinogen <100 mg/dL in DIC represents severe consumption and is a criterion for DIC diagnosis. In massive hemorrhage, fibrinogen <150--200 mg/dL is the threshold for fibrinogen replacement with cryoprecipitate or fibrinogen concentrate. If Clauss fibrinogen is significantly lower than fibrinogen antigen (immunologic assay), dysfibrinogenemia is suspected and further evaluation by thrombin time, reptilase time, and genetic testing may be pursued.

\section*{Understanding Results}

\subsection*{Below Normal (Hypofibrinogenemia: <200 mg/dL)}
\begin{itemize}
  \item \textbf{Disseminated intravascular coagulation (DIC):} Consumption of fibrinogen due to widespread thrombin generation and fibrin deposition. Fibrinogen is typically decreased (often <100 mg/dL) but may be ``normal'' if baseline was elevated (acute-phase reactant), masking consumption; a falling fibrinogen on serial measurements is more significant than a single value. Seen in sepsis, trauma, obstetric emergencies, malignancy, and severe pancreatitis.
  \item \textbf{Massive hemorrhage and transfusion:} Dilutional and consumptive coagulopathy. Fibrinogen is the first coagulation factor to reach critically low levels during massive transfusion protocol. Fibrinogen <150 mg/dL is an indication for cryoprecipitate or fibrinogen concentrate.
  \item \textbf{Congenital afibrinogenemia:} Autosomal recessive, absent fibrinogen, severe bleeding (umbilical stump, mucosal, intracranial, hemarthroses, post-surgical), diagnosed at birth. PT, aPTT, TT all markedly prolonged (no clot formation).
  \item \textbf{Congenital hypofibrinogenemia:} Autosomal recessive or dominant, low but detectable fibrinogen (20--100 mg/dL). Variable bleeding: usually mild or post-traumatic; women may have menorrhagia and pregnancy complications.
  \item \textbf{Advanced liver disease:} Impaired hepatic fibrinogen synthesis in cirrhosis, fulminant hepatic failure. Usually accompanied by prolonged PT and aPTT.
  \item \textbf{Thrombolytic therapy:} Administration of streptokinase, alteplase, or tenecteplase generates plasmin that degrades fibrinogen.
  \item \textbf{L-asparaginase therapy:} Chemotherapy for ALL depletes asparagine required for hepatic protein synthesis, reducing fibrinogen.
  \item \textbf{Severe malnutrition and malabsorption:} Reduced protein substrates for fibrinogen synthesis.
  \item \textbf{Hereditary fibrinogen storage disease:} Rare disorder with abnormal fibrinogen accumulation in hepatocytes.
\end{itemize}

\subsection*{Above Normal (Hyperfibrinogenemia: >400 mg/dL)}
\begin{itemize}
  \item \textbf{Acute and chronic inflammation:} Fibrinogen is a major acute-phase reactant. Elevated in infections (bacterial, viral), tissue injury, surgery, trauma, burns, and autoimmune diseases (rheumatoid arthritis, systemic lupus erythematosus, inflammatory bowel disease).
  \item \textbf{Malignancy:} Paraneoplastic elevation; particularly lung, pancreas, ovary, stomach, and kidney cancers. Fibrinogen promotes tumor angiogenesis, growth, and metastasis.
  \item \textbf{Cardiovascular disease:} Elevated fibrinogen is a well-established independent risk factor for coronary artery disease, myocardial infarction, stroke, and peripheral arterial disease. Promotes atherosclerosis, platelet aggregation, and thrombus formation.
  \item \textbf{Pregnancy:} Physiological elevation to 400--600 mg/dL in the third trimester; a prothrombotic adaptation to reduce postpartum hemorrhage risk.
  \item \textbf{Estrogen therapy:} Oral contraceptives and hormone replacement therapy increase fibrinogen concentration.
  \item \textbf{Diabetes mellitus:} Insulin resistance and hyperglycemia increase fibrinogen synthesis through chronic low-grade inflammation.
  \item \textbf{Smoking:} Dose-dependent increase in fibrinogen; levels decline after smoking cessation.
  \item \textbf{Advanced age:} Fibrinogen increases with age.
  \item \textbf{Obesity:} Adipose tissue-derived IL-6 stimulates hepatic fibrinogen synthesis.
  \item \textbf{Nephrotic syndrome:} Urinary loss of albumin and anticoagulant proteins (antithrombin) leads to compensatory hepatic synthesis of fibrinogen and other large proteins.
  \item \textbf{Genetic variation:} Polymorphisms in the fibrinogen beta-chain gene (G-455A) associated with higher fibrinogen levels.
\end{itemize}

\section*{Normal Results}
\begin{itemize}
  \item \textbf{Adults:} 200--400 mg/dL (2.0--4.0 g/L) by Clauss functional method
  \item \textbf{Newborns (first week):} 150--300 mg/dL (physiologically lower)
  \item \textbf{Infants (1--12 months):} 160--350 mg/dL
  \item \textbf{Pregnancy (third trimester):} 400--600 mg/dL (physiological elevation)
  \item \textbf{Critical values (action required):} <100 mg/dL (risk of spontaneous bleeding; fibrinogen replacement indicated); >700 mg/dL (increased thrombotic risk)
\end{itemize}

\section*{Follow-up Tests}
\begin{itemize}
  \item \textbf{Fibrinogen antigen assay (immunologic):} Measures total fibrinogen protein (functional plus non-functional). Discrepancy between Clauss (low) and antigen (normal) indicates dysfibrinogenemia.
  \item \textbf{Thrombin time (TT) and reptilase time:} Both prolonged in hypofibrinogenemia and dysfibrinogenemia. Reptilase time (uses batroxobin, not thrombin) is not affected by heparin.
  \item \textbf{PT and aPTT:} Both are prolonged in severe fibrinogen deficiency (no fibrinogen = no clot).
  \item \textbf{D-dimer:} To evaluate for DIC or fibrinolysis.
  \item \textbf{Complete blood count with platelet count:} For DIC evaluation (thrombocytopenia).
  \item \textbf{Liver function tests (ALT, AST, albumin, bilirubin, INR):} To assess hepatic synthetic function.
  \item \textbf{CRP (C-reactive protein):} To corroborate the acute-phase response in elevated fibrinogen.
  \item \textbf{Genetic testing (FGA, FGB, FGG genes):} For congenital fibrinogen disorders --- afibrinogenemia, hypofibrinogenemia, dysfibrinogenemia.
  \item \textbf{Thromboelastography (TEG) or rotational thromboelastometry (ROTEM):} Functional assessment of fibrinogen contribution to clot strength (FIBTEM assay) in massive hemorrhage.
  \item \textbf{Fibrinogen beta-chain gene polymorphism (G-455A):} When genetic contribution to elevated fibrinogen is suspected.
\end{itemize}

\section*{References}
\begin{enumerate}
  \item Tierney LM, Saint S, Drazen JM. CURRENT Medical Diagnosis \& Treatment. McGraw-Hill; 2023.
  \item Wallach J. Interpretation of Diagnostic Tests. 11th ed. Wolters Kluwer; 2022.
  \item Fischbach FT, Dunning MB. A Manual of Laboratory and Diagnostic Tests. 11th ed. Wolters Kluwer; 2023.
  \item Burtis CA, Ashwood ER, Bruns DE. Tietz Textbook of Clinical Chemistry and Molecular Diagnostics. 7th ed. Elsevier; 2023.
  \item Levi M, Ten Cate H. Disseminated intravascular coagulation. N Engl J Med. 1999;341(8):586--592.
  \item Taylor FB Jr, Toh CH, Hoots WK, et al. Towards definition, clinical and laboratory criteria, and a scoring system for DIC. Thromb Haemost. 2001;86(5):1327--1330.
  \item Danesh J, Lewington S, Thompson SG, et al. Plasma fibrinogen level and the risk of major cardiovascular diseases and nonvascular mortality. JAMA. 2005;294(14):1799--1809.
  \item de Moerloose P, Casini A, Neerman-Arbez M. Congenital fibrinogen disorders: an update. Semin Thromb Hemost. 2013;39(6):585--595.
\end{enumerate}

\clearpage
